\documentclass[../DoAn.tex]{subfiles}
\begin{document}

\begin{center}
    \Large{\textbf{TÓM TẮT NỘI DUNG BÁO CÁO}}\\
\end{center}
\vspace{1cm}
Hệ thống IoT và băng rộng hiện đại bao gồm nhiều thành phần phân tán như thiết bị CPE, router gia đình, ONT, nền tảng ACS/USP, hệ thống backend của nhà cung cấp dịch vụ và trung tâm SOC/SIEM. Các thành phần này thường có vòng đời dài, tài nguyên hạn chế, phụ thuộc firmware của nhiều nhà sản xuất và vận hành trên quy mô lớn, vì vậy rủi ro an ninh không thể được xử lý bằng một biện pháp đơn lẻ. Báo cáo nghiên cứu quy trình quản lý an ninh cho hệ thống IoT và băng rộng theo năm trụ cột: đánh giá lỗ hổng, mô hình hóa mối đe dọa, ghi nhật ký và kiểm toán, phản ứng sự cố, quản lý bản vá bảo mật. Nội dung được tổng hợp từ các chuẩn NIST SP 800-30, NIST SP 800-61, ISO/IEC 27035, MITRE ATT\&CK for IoT và danh mục bài báo khảo sát trong thư mục nguồn. Báo cáo xây dựng kiến trúc tham chiếu, sơ đồ luồng dữ liệu, quy trình đánh giá lỗ hổng, vòng đời phản ứng sự cố và quy trình OTA an toàn. Kết quả chính là một khung phân tích có thể dùng để nhận diện tài sản, đánh giá rủi ro, ưu tiên xử lý CVE/CVSS, thiết kế log phục vụ truy vết, tổ chức playbook phản ứng sự cố và triển khai vá lỗi theo vòng kiểm thử--ký số--canary--giám sát. Báo cáo cũng chỉ ra các khoảng trống nghiên cứu như thiếu benchmark firmware đáng tin cậy, khó thu thập log đồng nhất, hạn chế vá lỗi trên thiết bị legacy và nhu cầu tự động hóa điều phối giữa SOC, NOC và hệ thống quản trị thiết bị.
\begin{flushright}
Sinh viên thực hiện\\
\begin{tabular}{@{}c@{}}
\textit{(Ký và ghi rõ họ tên)}
\end{tabular}
\end{flushright}

\end{document}